\documentclass{beamer}

\usepackage[backend=biber,style=abnt,uniquename=init,giveninits,repeatfields]{biblatex} % Citações padrão ABNT
\addbibresource{ref.bib}
\usepackage[utf8]{inputenc} % Pacote necessário quando se utiliza pdflatex
\usepackage[T1]{fontenc} % Pacote necessário quando se utiliza pdflatex
\usepackage[brazil]{babel}
\resetcounteronoverlays{exx}
\usepackage{tikz}
\usetikzlibrary{trees}
\usepackage{blindtext}
\usepackage{enumerate}
\usepackage{longtable}
\usepackage{parskip}
\usepackage{color}
\usepackage{multirow}
\usepackage{graphicx}
\usepackage{booktabs}
\usepackage{fancybox}
\usetheme{Copenhagen}
\setbeamertemplate{frametitle}[default][center]
\setbeamertemplate{caption}[numbered]
\setlength{\parskip}{12pt}

\makeatother
\setbeamertemplate{footline}
{
  \leavevmode%
  \hbox{%
  \begin{beamercolorbox}[wd=.2\paperwidth,ht=2.25ex,dp=1ex,center]{author in head/foot}%
    \usebeamerfont{author in head/foot}\insertshortauthor
  \end{beamercolorbox}%
  \begin{beamercolorbox}[wd=.8\paperwidth,ht=2.25ex,dp=1ex,center]{title in head/foot}%
      \usebeamerfont{title in head/foot}\insertshorttitle
  \end{beamercolorbox}}%
  \vskip0pt%
}
\makeatletter

\title{Efeitos do Prejuízo no Sono na Funcionalidade e Cognição de Sujeitos com Transtornos de Humor}
\author[Bruno Montezano]{Bruno Braga Montezano}
\institute{Universidade Católica de Pelotas}

\titlegraphic{\includegraphics[width=2cm]{img/logoppgsc.png}\hspace*{4.75cm}%
    \includegraphics[width=2cm]{img/logoucpel.png}
}

\begin{document}

\begin{frame}

\maketitle

\end{frame}

\begin{frame}
\frametitle{Introdução}



\end{frame}

\begin{frame}



\end{frame}

\begin{frame}



\end{frame}

\begin{frame}
    \frametitle{Objetivo Geral}

    \centering
    \Large
    Avaliar o efeito da insônia/hipersonia na funcionalidade e
    cognição de sujeitos com transtornos de humor

\end{frame}

\begin{frame}
\frametitle{Métodos}

    \begin{block}{Delineamento}
        Coorte prospectivo
    \end{block}

    \begin{block}{Amostra}
        Adultos com idade entre 18 e 60 anos diagnosticados com TDM
    \end{block}

    \begin{block}{Primeira fase}
        Aconteceu entre 2012 e 2015
    \end{block}

    \begin{block}{Segunda fase}
        Reavaliação em 2017
    \end{block}

\end{frame}

\begin{frame}
\frametitle{Métodos}

    \begin{tikzpicture}[grow=0, level distance=6cm, sibling distance=25mm, <-]
        \node{Avaliação (n = 966)}
        child { node {Encaminhados por serviços de saúde (n = 230)}}
        child { node {Participantes de pesquisas (n = 134)}}
        child { node {Busca espontânea (n = 602)}}
    ;
    \end{tikzpicture}

    \centering
    \Ovalbox{585 sujeitos avaliados com TDM}

\end{frame}

\begin{frame}
\frametitle{Métodos}

    \begin{block}{Instrumentos}
        \setbeamerfont{footnote}{size=\tiny}
        \begin{itemize}
            \item \emph{Mini International Neuropsychiatric Interview} (MINI-PLUS)\footfullcite{sheehan_mini-international_1998}
        \end{itemize}
    \end{block}

\end{frame}

\begin{frame}
\frametitle{Métodos}

    \begin{block}{Processamento e análise de dados}
        \begin{itemize}
            \item \emph{Open Data Kit} (ODK)        
            \item R 4.0.2
            \begin{itemize}
                \item Análise univariada: 
                \item Análise bivariada: 
            \end{itemize}
        \end{itemize}
    \end{block}

\end{frame}

\begin{frame}
\frametitle{Aspectos Éticos}
    \Large
    \begin{itemize}
        \item Aprovado pelo Comitê de Ética em Pesquisa da UCPel (502.604)
        \item Todos os participantes assinaram o Termo de Consentimento Livre e Esclarecido antes de participarem do estudo
    \end{itemize}

\end{frame}

\begin{frame}
\frametitle{Resultados}



\end{frame} 
\begin{frame}

\end{frame}

\begin{frame}
\frametitle{Discussão}
\setbeamerfont{footnote}{size=\tiny}
    \begin{itemize}
        \item 
        \item
        \item
    \end{itemize}

\end{frame}

\begin{frame}
\frametitle{Discussão}

    \begin{itemize}
        \item
        \item
    \end{itemize}

\end{frame}

\begin{frame}
\frametitle{Conclusão}

    \begin{itemize}
        \item
        \item
    \end{itemize}

\end{frame}

\begin{frame}[shrink=20]
\frametitle{Referências}

\printbibliography

\end{frame}

\begin{frame}
\frametitle{Obrigado pela atenção!}

    \begin{block}{Email para contato:}
        bruno.montezano@sou.ucpel.edu.br
    \end{block}

\end{frame}

\end{document}
